% Human reconstruction; source ledger: ordinary_host_provenance.json.
% This is an inclusion fragment, with no preamble or document environment.
\subsection{The ordinary host and a fixed pair on a whole window}

Here $\mathbf1$ is the primitive character of conductor one.  The following
steps isolate the ordinary specialization without asserting the remaining
cohomological estimate.  All character bounds retain their centering at $s=1$.

\paragraph{OH1. The field.}
Let $q\equiv3\pmod8$ and $d\equiv7\pmod{16}$ be distinct primes, and put
\[
 T_e=\mathbb Q(\zeta_{2^{e+2}})^+,
 \qquad E_e=T_e(\sqrt{-q},\sqrt{-d})\quad(e\ge1).
\]
Then
\[
 \operatorname{Gal}(E_e/\mathbb Q)=C_{2^e}\times C_2\times C_2,
 \quad [E_e:\mathbb Q]=4\cdot2^e,
 \quad E_e^+=T_e(\sqrt{qd}).
\]
Indeed $T_e$ is cyclic of degree $2^e$ and ramifies only at $2$.
The three quadratic subfields of
$\mathbb Q(\sqrt{-q},\sqrt{-d})$ ramify at $q$ or $d$, so its intersection
with $T_e$ is $\mathbb Q$.  Complex conjugation is $(1,-1,-1)$ in the
displayed product.  In $\mathbb Q(\zeta_{2^{e+2}qd})$, an automorphism with
cyclotomic exponent $a$ fixes $E_e$ precisely when
\[
 a\equiv\pm1\pmod{2^{e+2}},\qquad
 J_q(a)=J_d(a)=1,\qquad J_r(a)=\Bigl(\frac ar\Bigr).
\]
The last assertion follows either from quadratic Gauss sums or the quadratic
subfield theorem for cyclotomic fields: the discriminants of
$\mathbb Q(\sqrt{-q})$ and $\mathbb Q(\sqrt{-d})$ are $-q$ and $-d$.
Thus this is exactly the field specified by \texttt{IsChiField} for
$\mathbf1$.  Its conductor is $2^{e+2}qd$.
The source's order parameter is $n=2^e$, not the field degree $4n$:
\texttt{ordAny} gives $1\le c\le\operatorname{ord}(\mathbf1)=1$, so its
general formula $n=c2^e$ has $c=1$.

\paragraph{OH2. The two inputs, with their quantifiers.}
Fix $0<\varepsilon<1/52$ and define
\[
 e(x)=\left\lceil(2+8\varepsilon)\log_2x\right\rceil\quad(x\ge1),
 \qquad x^{2+8\varepsilon}\le2^{e(x)}\le2x^{2+8\varepsilon}.
\]
Write $\operatorname{Good}(\varepsilon,Y,X;q,d)$ for the complete source
condition \texttt{PairGood 1}: besides its zero-free boxes, it includes
\[
 \tfrac12Y^{\varepsilon/5}<\log q\le Y^{\varepsilon/5},\quad
 \tfrac12Y^{\varepsilon/5}(\log Y)^2<\log d
       \le Y^{\varepsilon/5}(\log Y)^2,
\]
$q\mid d+1$, and $p\mid d+1$ for every odd prime $p\le Y^{\varepsilon/5}$.
Its three nontrivial twists $\chi=J_q,J_d,J_qJ_d$ satisfy
$|G_\chi(u)|\le u^{1/2+\varepsilon/2}$ and
$|L'/L(1,\chi)|\le4\log u$ for every $u\in[Y,X]$.
The pair-selection input says that for every fixed $A\ge1$ this condition
holds for some pair at every sufficiently large $Y$, with $X=Y^A$.

For a number field $E$, let $S_p$ consist of its places above $p$ and its
infinite places, and let $G_{E,S_p}$ be the Galois group of its maximal
extension unramified outside $S_p$. Put
$D_{p,j}=(\mathbb Q_p/\mathbb Z_p)(-j)$, the usual Tate-twisted module, and define
\[
\ell_p(E,j)=\operatorname{length}_{\mathbb Z_p}\ker\left(
H^1(G_{E,S_p},D_{p,j})\longrightarrow\prod_{w\mid p}H^1(E_w,D_{p,j})\right).
\]
The cohomological size is
\[
 \mathcal V(E,t)=\sup_{S\text{ finite primes}}
       \sum_{p\in S}\log p\sum_{1\le j<t}\ell_p(E,j).
\]
The cohomological input also asserts finiteness of every displayed length.
The specialization of \texttt{VWinAt} for $\mathbf1$ and $401$ used here says:
for this choice of $e(x)$, there are $J\ge1,X_0$ such that, for every
$X\ge X_0$ and every pair satisfying the preceding congruences and size
conditions at $Y=X^{1/J}$, there is $W$ with
\[
 X^{1/J}\le W,\quad W^2\le X,\quad
 \mathcal V(E_{e(x)},\lfloor x^{401}\rfloor)
       \le \exp(\tfrac12x^\varepsilon)\quad(W\le x\le W^2).
\]
This input is the unexpanded arithmetic theorem in this subsection.

\paragraph{OH3. Fixed-pair assembly.}
Assuming the two inputs, there is $J\ge1$ such that for every lower
threshold $B$ there are $W\ge B$, $Y$, and a \emph{single} pair $q,d$ for
which, simultaneously for every $x\in[W,W^2]$,
\[
 \begin{gathered}
 x^{1/J}\le Y\le x,\qquad \operatorname{Good}(\varepsilon,Y,x;q,d),\\
 x^{2+8\varepsilon}\le n(x)=2^{e(x)}\le2x^{2+8\varepsilon},\qquad
 \mathcal V(E_{e(x)},\lfloor x^{400}\rfloor)
       \le\exp(\tfrac12x^\varepsilon).
 \end{gathered}
\]
To prove this, first fix $J,X_0$ from the cohomological input.  Apply pair
selection with $A=J$.  Choose $X\ge X_0$ so large that
$Y=X^{1/J}\ge\max(1,B)$ and exceeds its selection threshold.  Choose
$q,d$ good on $[Y,Y^J]=[Y,X]$; then choose $W$ from the cohomological
input for this very pair.  For $W\le x\le W^2\le X$, restrict the good
pair's upper endpoint to $x$ and decrease the twist cutoff from
$\lfloor x^{401}\rfloor$ to $\lfloor x^{400}\rfloor$.
All summands are nonnegative, so the bound persists.  Only $e(x)$ changes
with $x$.  The general source proof rescales by $c^{1/(2+8\varepsilon)}$
and retains $[W,W^{3/2}]$; $c=1$ makes that loss unnecessary here.

\paragraph{OH4. The four characters and the sign.}
The characters trivial on $C_{2^e}$ are exactly
$1,J_q,J_d,J_qJ_d$.  Weighting them by $+,+,-,-$ gives
\[
 S_{q,d}(x)=G_{\mathbf1}(x)+G_{J_q}(x)-G_{J_d}(x)-G_{J_qJ_d}(x),
 \qquad |S_{q,d}(x)-G_{\mathbf1}(x)|\le3x^{1/2+\varepsilon/2}.
\]
At a prime power $p^k$, the corresponding von Mangoldt coefficient,
in units of $\log p$, is
$(1+J_q(p)^k)(1-J_d(p)^k)$.
For $p\nmid qd$, it is $4$ exactly when $k$ is odd,
$J_q(p)=1$, and $J_d(p)=-1$; otherwise it is zero.
Since $d\equiv-1\pmod q$, reciprocity gives $J_q(d)=-1$ and
$J_d(q)=1$.  Thus all terms at $q$ vanish; at $d$, the even powers have
coefficient $2$ and the odd powers vanish.  Also $J_d(2)=1$, and the
small-prime congruences give $J_d(p)=1$ for all odd
$p\le Y^{\varepsilon/5}$.  This identifies the prime class exactly.
It does not replace the centered ramp bounds, zero-free boxes, or the
cohomological input by a bound on a finite uncentered prime sum.
